\documentclass[12pt,a4paper]{article}
\usepackage{amsmath,amssymb,mathrsfs,tikz,times,pifont}
\usepackage{enumitem}
\usepackage{multicol}
\usepackage{lmodern}
\newcommand\circitem[1]{%
\tikz[baseline=(char.base)]{
\node[circle,draw=gray, fill=red!55,
minimum size=1.2em,inner sep=0] (char) {#1};}}
\newcommand\boxitem[1]{%
\tikz[baseline=(char.base)]{
\node[fill=cyan,
minimum size=1.2em,inner sep=0] (char) {#1};}}
\setlist[enumerate,1]{label=\protect\circitem{\arabic*}}
\setlist[enumerate,2]{label=\protect\boxitem{\alph*}}
\everymath{\displaystyle}
\usepackage[left=1cm,right=1cm,top=1cm,bottom=1.7cm]{geometry}
\usepackage[colorlinks=true, linkcolor=blue, urlcolor=blue, citecolor=blue]{hyperref}
\usepackage{array,multirow}
\usepackage[most]{tcolorbox}
\usepackage{varwidth}
\usepackage{float}
\tcbuselibrary{skins,hooks}
\usetikzlibrary{patterns}

\newtcolorbox{exa}[2][]{enhanced,breakable,before skip=2mm,after skip=5mm,
colback=yellow!20!white,colframe=black!20!blue,boxrule=0.5mm,
attach boxed title to top left ={xshift=0.6cm,yshift*=1mm-\tcboxedtitleheight},
fonttitle=\bfseries,
title={#2},#1,
boxed title style={frame code={
\path[fill=tcbcolback!30!black]
([yshift=-1mm,xshift=-1mm]frame.north west)
arc[start angle=0,end angle=180,radius=1mm]
([yshift=-1mm,xshift=1mm]frame.north east)
arc[start angle=180,end angle=0,radius=1mm];
\path[left color=tcbcolback!60!black,right color = tcbcolback!60!black,
middle color = tcbcolback!80!black]
([xshift=-2mm]frame.north west) -- ([xshift=2mm]frame.north east)
[rounded corners=1mm]-- ([xshift=1mm,yshift=-1mm]frame.north east)
-- (frame.south east) -- (frame.south west)
-- ([xshift=-1mm,yshift=-1mm]frame.north west)
[sharp corners]-- cycle;
},interior engine=empty,
},interior style={top color=yellow!5}}

\usepackage{fancyhdr}
\usepackage{eso-pic}
\usepackage{tkz-tab}
\AddToShipoutPicture{
    \AtTextCenter{%
        \makebox[0pt]{\rotatebox{80}{\textcolor[gray]{0.7}{\fontsize{5cm}{5cm}\selectfont PGB}}}
    }
}
\usepackage{lastpage}
\fancyhf{}
\pagestyle{fancy}
\renewcommand{\footrulewidth}{1pt}
\renewcommand{\headrulewidth}{0pt}
\renewcommand{\footruleskip}{10pt}
\fancyfoot[R]{\color{blue}\ding{45}\ \textbf{2025}}
\fancyfoot[L]{\color{blue}\ding{45}\ \textbf{Prof : M. BA}}
\cfoot{\bf \thepage / \pageref{LastPage}}

%----------------------------------
% Compteurs
%----------------------------------

\newcounter{probleme}
\newcounter{annee}
\newcounter{session}

\setcounter{probleme}{0}
\setcounter{annee}{1999}   % année de départ
\setcounter{session}{0}    % 0 = 1er groupe, 1 = Remplacement

%----------------------------------
% Commande automatique
%----------------------------------

\newcommand{\exo}{%
    \refstepcounter{probleme}%
    \par\medskip
    \noindent\textbf{\underline{Problème \theprobleme}\quad
    BAC \theannee\ %
    \ifnum\value{session}=0
        1\textsuperscript{er} GROUPE%
        \setcounter{session}{1}%
    \else
        Remplacement%
        \setcounter{session}{0}%
        \stepcounter{annee}%
    \fi
    }
    \par\medskip
}


\begin{document}
\renewcommand{\arraystretch}{1.5}
\renewcommand{\arrayrulewidth}{1.2pt}
\begin{tikzpicture}[overlay,remember picture]
    \node[draw=blue,line width=1.2pt,fill=purple,text=blue,inner sep=3mm,rounded corners,pattern=dots]at ([yshift=-2.5cm]current page.north) {\begingroup\setlength{\fboxsep}{0pt}\colorbox{white}{\begin{tabular}{|*1{>{\centering \arraybackslash}p{0.28\textwidth}} |*2{>{\centering \arraybackslash}p{0.2\textwidth}|} *1{>{\centering \arraybackslash}p{0.19\textwidth}|} }
                \hline
                \multicolumn{3}{|c|}{$\diamond$$\diamond$$\diamond$\ \textbf{Lycée de Dindéfélo}\ $\diamond$$\diamond$$\diamond$ } & \textbf{A.S. : 2024/2025} \\ \hline
                \textbf{Matière : Mathématiques} & \textbf{Niveau : T S2} & \textbf{Date : 09/06/2025} & \textbf{} \\ \hline
                \multicolumn{4}{|c|}{\parbox[c]{10cm}{\begin{center}
                  \textbf{{\Large\sffamily Problèmes proposés au BAC S2 Sénégal de 1999 à 2022}}
                \end{center}}} \\ \hline
            \end{tabular}}\endgroup};
\end{tikzpicture}
\vspace{3cm}

\begin{multicols}{2}
\small
\setlength{\columnseprule}{0.1mm}

\exo

On considère la fonction \( f \) définie par :

\( 
f(x) =
\begin{cases}
x + \ln\left|\dfrac{x - 1}{x + 1}\right| & \text{si } x \in ]-\infty, -1[ \cup ]-1, 0[ \\
x^2 e^{-x} & \text{si } x \in [0, +\infty[
\end{cases}
\)
et \( (C_f) \) sa courbe représentative dans un repère orthonormé \( (O; \vec{i}, \vec{j}) \), d’unité 2 cm.

\textbf{Partie A :} 
\begin{enumerate}
  \item Déterminer l’ensemble de définition \( D_f \) de \( f \). Calculer \( f(-2) \) et \( f(3) \).
  \item Calculer les limites aux bornes de \( D_f \).
  \item Étudier la continuité de \( f \) en 0.
  \item 
  \begin{enumerate}
    \item Établir que la dérivée de \( f \) est donnée par :
    
\(
    f'(x) =
    \begin{cases}
    \dfrac{x^2 + 1}{x^2 - 1}  \text{si } x \in ]-\infty, -1[ \cup ]-1, 0[ \\
    x e^{-x}(2 - x) \text{si } x \in [0, +\infty[
    \end{cases}
\)
    \item La fonction \( f \) est-elle dérivable en 0 ? Justifier votre réponse.
    \item Dresser le tableau de variations de \( f \).
  \end{enumerate}
  \item Démontrer que l’équation \( f(x) = 0 \) admet une solution unique \( \alpha \) comprise entre \( -1{,}6 \) et \( -1{,}5 \).
  \item 
  \begin{enumerate}
    \item Justifier que la droite \( (D) \) d’équation 
    
    \( y = x \) est une asymptote à la courbe \( (C_f) \) en \( -\infty \).
    \item Étudier la position relative de \( (C_f) \) par rapport à la droite \( (D) \) pour
    
     \( x \in ]-\infty, -1[ \cup ]-1, 0[ \).
    \item Tracer \( (C_f) \).
  \end{enumerate}
\end{enumerate}

\textbf{Partie B :}  Soit \( g \) la restriction de \( f \) à \( I = [0;2] \).
\begin{enumerate}
    \item Montrer que \( g \) définit une bijection de \( I \) vers un intervalle \( J \) à préciser.
    \item On note \( g^{-1} \) la bijection réciproque de \( g \).
    \begin{enumerate}
        \item Résoudre l’équation \( g^{-1}(x) = 1 \).
        \item Montrer que \( \left(g^{-1}\right)'\left(\dfrac{1}{e}\right) = e \).
        \item Construire \( (C_{g^{-1}}) \), la courbe de \( g^{-1} \).
    \end{enumerate}
\end{enumerate}


\textbf{Partie C :}  \( \beta \) étant un réel strictement positif, on pose :

\(I(\beta) = \int_0^{\beta} f(x)\,dx \)
\begin{enumerate}
    \item 
    \begin{enumerate}
        \item Interpréter graphiquement \( I(\beta) \).
        \item En procédant par une intégration par parties, calculer \( I(\beta) \).
    \end{enumerate}
    
    \item Calculer \( \lim_{\beta \to -\infty} I(\beta) \).
    
    \item On pose \( \beta = 2 \).
    \begin{enumerate}
        \item Calculer \( I(2) \).
        \item En déduire la valeur en \( \text{cm}^2 \) de l’aire du domaine du plan délimité par la courbe \( (C_f) \), l’axe des abscisses et les droites d’équations \( x = 0 \) et \( x = \dfrac{4}{e^2} \).
    \end{enumerate}
\end{enumerate}

\exo
Soit \( f \) la fonction définie sur \( \mathbb{R} \) par :\\
\(
f(x) =
\begin{cases}
e^{-\frac{1}{x^2}} & \text{si } x \in ]-\infty ; 0[ \\
\ln\left|\dfrac{x - 1}{x + 1}\right| & \text{si } x \in [0 ; 1[ \cup ]1 ; +\infty[
\end{cases}
\)\\
et \( (C_f) \) sa courbe représentative dans un repère orthonormé \( (O ; \vec{i} ; \vec{j}) \) d’unité \( 2\ cm \).

\textbf{Partie A :} 
\begin{enumerate}
    \item Étudier la continuité de \( f \) en 0.
    \item 
    \begin{enumerate}
        \item Montrer que \( \forall x \in ]0 ; 1[ \),\\
        \( \dfrac{f(x)}{x} = \dfrac{\ln(1 - x)}{x} - \dfrac{\ln(1 + x)}{x} \)
        \item Étudier la dérivabilité de \( f \) en 0.
        \item En déduire que \( (C_f) \) admet au point d’abscisse 0 deux demi-tangentes dont on donnera les équations.
    \end{enumerate}
        \item Étudier les variations de \( f \).
				\item Tracer \( (C_f) \).
\end{enumerate}				
    \textbf{Partie B :}  Soit \( g \) la restriction de \( f \) à \( ]1 ; +\infty[ \).
\begin{enumerate}

    \item Montrer que \( g \) est une bijection de \( ]1 ; +\infty[ \) vers un intervalle \( J \) à préciser.\\
    On notera \( g^{-1} \) la bijection réciproque de \( g \).

    \item Montrer que l’équation \( g(x) = -e \) admet une unique solution \( \alpha \) sur l’intervalle \( ]1 ; +\infty[ \).\\
    \textbf{(On ne demande pas de calculer \( \alpha \)).}

    \item Montrer que \( \forall x \in J \), \( g^{-1}(x) = 1 - \dfrac{e^x}{e^x - 1} \).

    \item Construire \( (C_{g^{-1}}) \).\\
    \textbf{(On indiquera la nature et l’équation de chacune des asymptotes à \( (C_g) \) et \( (C_{g^{-1}}) \)).}
\item Calculer en \( \text{cm}^2 \) l’aire \( A \) de l’ensemble des points \( M(x ; y) \) défini par :

\(
\left\{
\begin{array}{l}
-\ln 7 \leq x \leq -1 \\
0 \leq y \leq g^{-1}(x)
\end{array}
\right.
\)

\end{enumerate}

\exo
Soit \( f \) la fonction définie de \( \mathbb{R} \) dans \( \mathbb{R} \) par :\\
\(
f(x) =
\begin{cases}
x \ln(x+1) & \text{si } x \geq 0 \\
x e^{\frac{1}{x}} & \text{si } x < 0
\end{cases}
\)

Le plan est muni d’un repère orthonormé \( (O ; \vec{i} ; \vec{j}) \) \textbf{(unité graphique 2 cm)}.\\
On désigne par \( (C) \) la courbe représentative de \( f \) et \( (\Delta) \) la droite d’équation \( y = x \).

\textbf{Partie A :}
\begin{enumerate}
    \item 
    \begin{enumerate}
        \item Montrer que \( f \) est continue en 0.
        \item Étudier la dérivabilité de \( f \) en 0.
        \item Interpréter les résultats précédents.
    \end{enumerate}

    \item 
    \begin{enumerate}
        \item Montrer que \( \forall x < 0 \), \( f'(x) > 0 \).
        \item Étudier les variations de \( f \) sur \( [0 ; +\infty[ \).\\
        En déduire que \( \forall x > 0 \), \( f'(x) > 0 \).
        \item Donner le tableau de variations de \( f \).
    \end{enumerate}
    \item
\begin{enumerate}
    \item Déterminer \( \lim\limits_{x \to -\infty} x \left( e^{\frac{1}{x}} - 1 \right) \), \textbf{(on pourra poser \( u = \dfrac{1}{x} \))}.
    \item Montrer que la droite \( (D) : y = x + 1 \) est asymptote à \( (C) \) au voisinage de \( -\infty \).\\
    On admettra que \( (C) \) est en dessous de \( (D) \).
    \item Déterminer la nature de la branche infinie de \( (C) \) en \( +\infty \).
\end{enumerate}

\item Construire \( (C) \), on précisera les coordonnées de \( I \), intersection de \( (C) \) et \( (\Delta) \) pour \( x > 0 \).

\end{enumerate}
\textbf{Partie B :}
\begin{enumerate}
    \item Déterminer les réels \( a \), \( b \) et \( c \) tels que :
    
    \(
    \forall x \in \mathbb{R}_+, \quad \dfrac{x^2}{x + 1} = ax + b + \dfrac{c}{x + 1}.
    \)

    \item En déduire au moyen d’une intégration par parties que la fonction \( F \) telle que :

    \(
    F(x) = \dfrac{(x^2 - 1)\ln(x + 1)}{2} - \dfrac{1}{4}(x^2 - 2x)
    \)

    est une primitive de \( f \) sur \( \mathbb{R}_+ \).

    \item En déduire en \( \text{cm}^2 \) l’aire \( A \) de la partie du plan délimitée par \( (\Delta) \), \( (C) \) et les droites d’équations \( x = 0 \) et \( x = e - 1 \).
\end{enumerate}

\textbf{Partie C :}
\begin{enumerate}
\item 
    \begin{enumerate}
        \item Montrer que \( f \) admet une bijection réciproque notée \( f^{-1} \).
        \item \( f^{-1} \) est-elle dérivable en 0 ? Préciser la nature de la tangente en 0 à la courbe \( f^{-1} \).
    \end{enumerate}
    \item Construire \( (C_0) \), la courbe de \( f^{-1} \) dans le repère \( (O ; \vec{i} ; \vec{j}) \).

    \item Déduire du \textbf{B.3)} l’aire du domaine \( D \) définie par : \( M(x ; y) \) tels que :
    
    \(
    \left\{
    \begin{array}{l}
    0 \leq x \leq e - 1 \\
    f(x) \leq y \leq f^{-1}(x)
    \end{array}
    \right.
    \)
\end{enumerate}

\exo

\textbf{Partie A :} Soit \( g \) la fonction définie par 

\( g(x) = 1 - x e^{-x} \).

\begin{enumerate}
    \item Étudier les variations de \( g \).
    \item En déduire le signe de \( g(x) \) suivant les valeurs de \( x \).
\end{enumerate}

\textbf{Partie B :} Soit \( f \) la fonction définie par :\\
\(
f(x) =
\begin{cases}
\ln(-x) & \text{si } x < -1 \\
(x + 1)(1 + e^{-x}) & \text{si } x \leq -1
\end{cases}
\)

On désigne par \( (C) \) sa courbe représentative dans un repère orthonormé \( (0 ; \vec{i} ; \vec{j}) \) \textbf{(unité graphique 2 cm)}.

\begin{enumerate}
    \item Étudier la continuité et la dérivabilité de \( f \) sur \( \mathbb{R} \).

    \item Étudier les variations de \( f \), puis dresser le tableau de variations de \( f \).

    \item
    \begin{enumerate}
        \item Montrer que la droite \( (D) : y = x + 1 \) est une asymptote à \( (C_f) \) en \( +\infty \).
        \item Étudier la position relative de \( (C_f) \) par rapport à \( (D) \) sur \( [-1 ; +\infty[ \).
    \end{enumerate}

    \item Montrer qu’il existe un unique point de la courbe \( (C_f) \) où la tangente \( (T) \) est parallèle à la droite \( (D) \).

    \item Tracer \( (C_f) \), l’asymptote \( (D) \) et la tangente \( (T) \),\\
    on précisera la tangente ou les demi-tangentes à \( (C_f) \) au point d’abscisse \( -1 \).

    \item
    \begin{enumerate}
        \item Montrer que \( f \) est une bijection de \( [-1 ; +\infty[ \) sur un ensemble \( J \) que l’on précisera.
        \item Construire \( (C_0) \), la courbe de \( f^{-1} \) sur le même graphique que la courbe \( (C_f) \).
    \end{enumerate}
\end{enumerate}

\textbf{Partie C :} Pour \( \beta \geq -1 \), on note \( A(\beta) \) l’aire en \( \text{cm}^2 \) de la partie du plan définie par :

\(
\left\{
\begin{array}{l}
-1 \leq x \leq \beta \\
x + 1 \leq y \leq f(x)
\end{array}
\right.
\)

\begin{enumerate}
    \item Calculer \( A(\beta) \) à l’aide d’une intégration par parties.
    \item Montrer que \( A(\beta) \) admet une limite finie lorsque \( \beta \to +\infty \).\\
    Interpréter graphiquement cette limite.
\end{enumerate}

\exo

On considère la fonction \( g \) définie par :\\
\(
f(x) =
\begin{cases}
x(1 - \ln x)^2 & \text{si } x > 0 \\
0 & \text{si } x = 0
\end{cases}
\)

On appelle \( (C) \) sa courbe représentative dans un repère orthonormé \( (O ; \vec{i} ; \vec{j}) \).

\begin{enumerate}
    \item Étudier la continuité et la dérivabilité de \( g \) sur son ensemble de définition.
    \item Étudier les variations de \( g \) puis dresser son tableau de variations.
    \item Tracer \( (C) \).
    \item Soit \( \alpha \) un réel appartenant à l’intervalle \( ]0 ; e[ \).
    \begin{enumerate}
        \item Calculer à l’aide de deux intégrales par parties, l’aire \( A(\alpha) \) du domaine plan délimitée par l’axe des abscisses, la courbe \( (C) \) et les droites d’équations respectives \( x = \alpha \) et \( x = e \).
        \item Calculer
    \end{enumerate}
    \item
    \begin{enumerate}
        \item Déterminer les coordonnées des points d’intersection de la courbe \( (C) \) et la droite \( (\Delta) \) d’équation \( y = x \).
        \item Pour quelles valeurs de \( m \) la droite \( (\Delta_m) \) d’équation \( y = mx \) recoupe-t-elle la courbe \( (C) \) en deux points \( M_1 \) et \( M_2 \) autres que \( O \) ?
        \item La droite \( (\Delta_m) \) coupe la droite \( (D) \) d’équation \( x = e \) en \( P \). 
        
        Montrer que \( OM_1 \times OM_2 = OP^2 \).
    \end{enumerate}

    \item
    \begin{enumerate}
        \item Montrer que la restriction \( h \) de la fonction \( g \) à l’intervalle \( [e ; +\infty[ \) admet une réciproque \( h^{-1} \) dont on précisera l’ensemble de définition.
        \item Sur quel ensemble \( h^{-1} \) est-elle dérivable ?\\
        Calculer \( h(e^2) \) ; en déduire \( (h^{-1})'(e^2) \).
        \item Construire la courbe de \( h^{-1} \).
    \end{enumerate}

\end{enumerate}

\exo
\textbf{vide}
\exo

\textbf{Partie A :} On considère la fonction \( g \) définie sur \( \mathbb{R}^+ \setminus \{1\} \) par :

\(
g(x) =
\begin{cases}
\dfrac{1}{(\ln x)^2} - \dfrac{1}{\ln x} & \text{si } x > 0 \text{ et } x \neq 1 \\
0 & \text{si } x = 0
\end{cases}
\)

\begin{enumerate}
    \item Montrer que \( g \) est continue en 0.
    \item Étudier les limites de \( g \) aux bornes de son ensemble de définition.
    \item Dresser le tableau de variations de \( g \).
    \item En déduire le signe de \( g(x) \) en fonction de \( x \).
    \item Calculer en \( \text{cm}^2 \) l’aire de la partie plane comprise entre la courbe de \( g \), l’axe des abscisses et les droites d’équations respectives : \( x = e \) et \( x = e^2 \).
\end{enumerate}

\textbf{Partie B :} On considère la fonction \( f \) définie sur \( \mathbb{R}^+ \setminus \{1\} \) par :
\(
f(x) =
\begin{cases}
-\dfrac{x}{\ln x} & \text{si } x > 0 \text{ et } x \neq 1 \\
0 & \text{si } x = 0
\end{cases}
\)

\begin{enumerate}
    \item Montrer que \( f \) est continue à droite et dérivable à droite au point 0.\\
    En déduire l’existence d’une demi-tangente à la courbe représentative \( (C) \) de \( f \) au point d’abscisse 0.

    \item Étudier les limites aux bornes de son ensemble de définition.

    \item Comparer \( f'(x) \) et \( g(x) \).\\
    En déduire les variations de \( f \) et son tableau de variations.

    \item Déterminer l’équation de la tangente \( (D) \) à la courbe \( (C) \) au point d’abscisse \( e^2 \).
    \item Soit \( M \) le point de \( (C) \) d’abscisse \( x \) et \( N \) le point de \( (D) \) de même abscisse \( x \).\\
    On pose \( \varphi(x) = \overline{MN} \).
    \begin{enumerate}
        \item Montrer que \( \varphi(x) = f(x) + \dfrac{x + e^2}{4} \).
        \item Déduire de la \textbf{partie A} le tableau de variations de \( f'(x) \) puis le signe de \( f'(x) \) sur \( ]1 ; +\infty[ \).
        \item En déduire le signe de \( \varphi(x) \) sur \( ]1 ; +\infty[ \) et la position de \( (C) \) par rapport à \( (D) \) pour les points d’abscisse \( x > 1 \).
    \end{enumerate}

    \item Représenter dans le plan rapporté à un repère orthonormé la courbe \( (C) \) et la droite \( (D) \) \textbf{unité 2cm}.

\end{enumerate}
\exo
\textbf{vide}
\exo

\textbf{Partie A :} On considère la fonction \( u \) définie sur \( [0 ; +\infty[ \) par :
\[
u(x) = \ln\left| \frac{x+1}{x-1} \right| - \frac{2x}{x^2 - 1}
\]

\begin{enumerate}
    \item 
    \begin{enumerate}
        \item Déterminer l’ensemble de définition de \( u \).
        \item Calculer \( u(0) \) et \( \lim\limits_{x \to +\infty} u(x) \).
    \end{enumerate}

    \item Étudier les variations de \( u \).\\
    \textbf{(il n’est pas nécessaire de calculer la limite de \( u \) en 1).}

    \item Déduire des résultats précédents que :
    \begin{enumerate}
        \item \( \forall x \in [0 ; 1[, \ u(x) \geq 0 \).
        \item \( \forall x \in ]1 ; +\infty[, \ u(x) < 0 \).
    \end{enumerate}
\end{enumerate}

\textbf{Partie B :} Soit \( g \) la fonction définie sur \( [0 ; +\infty[ \) par :
\[
g(x) = x \ln\left| \frac{x+1}{x-1} \right| - 1
\]

\begin{enumerate}
    \item Déterminer \( D_g \) (le domaine de définition de \( g \)) ;\\
    puis étudier la limite de \( g \) en 1.
\end{enumerate}
\begin{enumerate}
    \item
    \begin{enumerate}
        \item Vérifier que \( \dfrac{x+1}{x-1} = 1 + \dfrac{2}{x-1} \).
        \item En déduire que 
        \[
        \lim_{x \to +\infty} \dfrac{x - 1}{2} \ln\left(1 + \dfrac{2}{x - 1} \right) = 1.
        \]
        \item En déduire que \( \lim\limits_{x \to +\infty} g(x) = 1 \). Interpréter géométriquement ce résultat.
        \item Dresser le tableau de variations de \( g \).
        \item Montrer qu’il existe un réel \( \alpha \) unique appartenant à \( ]0 ; 1[ \) tel que \( g(\alpha) = 0 \).\\
        Donner un encadrement d’ordre 1 de \( \alpha \).
    \end{enumerate}

    \item Tracer la courbe \( (C_g) \) de \( g \) dans le plan rapporté à un repère orthonormé \textbf{(unité 2 cm)}.
\end{enumerate}

\textbf{Partie C :} Soit \( h \) la fonction définie sur \( [0 ; 1[ \) et
\[
f(x) = (x^2 - 1) \ln \sqrt{ \dfrac{x+1}{x-1} }.
\]

\begin{enumerate}
    \item Montrer que \( f \) est dérivable sur \( [0 ; 1[ \) et que :\\
    \( f'(x) = g(x) \), \( \forall x \in [0 ; 1[ \).
    
    \item Déterminer l’aire du domaine plan limité par la courbe \( (C_g) \), l’axe des abscisses et la droite d’équation \( x = \alpha \).
\end{enumerate}
\exo
\textbf{vide}
\exo

Soit \( f \) la fonction définie par :
\[
f(x) = \frac{(2x - 1)e^x - 2x + 2}{e^x - 1}.
\]

On note \( (C) \) la représentation graphique de la fonction \( f \) dans un repère orthonormé \( (O ; \vec{i} ; \vec{j}) \), dont l’unité est 2 cm.

\begin{enumerate}
    \item Déterminer l’ensemble de définition \( D_f \) de la fonction \( f \), et trouver les réels \( a, b \) et \( c \) tels que pour tout \( x \in D_f \), on ait :
    \[
    f(x) = ax + b + \frac{c}{e^x - 1}.
    \]

    \item Déterminer les limites de \( f \) aux bornes de \( D_f \).

    \item 
    \begin{enumerate}
        \item Déterminer la fonction dérivée de \( f \).
        \item Résoudre dans \( \mathbb{R} \), l’équation :
        \[
        2e^{2x} - 5e^x + 2 = 0.
        \]
        \item En déduire le sens de variation de \( f \) et dresser le tableau de variations de \( f \).
    \end{enumerate}

    \item Démontrer que les droites d’équations respectives \( y = 2x - 1 \) et \( y = 2x - 2 \) sont des asymptotes de \( (C) \) respectivement en \( +\infty \) et en \( -\infty \).\\
    Préciser l’autre asymptote.

    \item Soit \( x \) un réel de \( D_f \). On considère les deux points \( M \) et \( M' \) de \( (C) \) d’abscisses respectives \( x \) et \( -x \).\\
    Déterminer les coordonnées du milieu \( \Omega \) du segment \( [MM'] \).\\
    Que peut-on en déduire pour la courbe \( (C) \) ?
    \item Tracer la courbe \( (C) \).
    \item
    \begin{enumerate}
        \item Trouver les réels \( \alpha \) et \( \beta \) tels que, pour tout réel \( x \) de l’ensemble \( D_f \), on ait :
        \[
        f(x) = 2x + \alpha + \frac{\beta e^x}{e^x - 1}.
        \]

        \item Soit \( k \) un réel supérieur ou égal à 2.\\
        Déterminer l’aire \( A(k) \) en \( \text{cm}^2 \) de l’ensemble des points du plan dont les coordonnées \( (x ; y) \) vérifient :
        \[
        \ln 2 \leq x \leq \ln k \quad \text{et} \quad 2x - 1 \leq y \leq f(x).
        \]
\end{enumerate}

\end{enumerate}

\exo

\textbf{Partie A :} Soit l’équation différentielle
\[
(E) : -\tfrac{1}{2}y'' + \tfrac{3}{2}y' - y = 0
\]
Déterminer la solution \( g \) de \( (E) \) dont la courbe représentative \( (C) \) passe par le point \( A(0 ; -1) \) et dont la tangente en ce point est parallèle à l’axe des abscisses.

\textbf{Partie B :} Soit la fonction définie sur \( \mathbb{R} \) par :
\[
f(x) = e^{2x} - e^x.
\]
On note \( (\Gamma) \) sa courbe représentative dans un repère orthonormé \( (O ; \vec{i} ; \vec{j}) \) \textbf{(unité 2 cm)}.

\begin{enumerate}
    \item Étudier les variations de \( f \).
    \item Déterminer l’équation de la tangente à \( (\Gamma) \) au point d’abscisse \( \ln 2 \).
    \item Calculer \( \lim\limits_{x \to +\infty} \dfrac{f(x)}{x} \). Interpréter graphiquement le résultat.
\end{enumerate}

\textbf{Partie C :} Soit \( h \) la restriction de \( f \) à l’intervalle \( [0 ; +\infty[ \).

\begin{enumerate}
    \item Démontrer que \( h \) réalise une bijection de \( [0 ; +\infty[ \) sur un intervalle \( J \) à préciser.
    \item Démontrer que \( h^{-1} \) est dérivable en \( 3 \), puis calculer \( (h^{-1})'(3) \).
    \item Déterminer \( h^{-1}(x) \) pour tout \( x \in J \).
    \item Tracer \( (C_0) \), la courbe représentative de \( h^{-1} \) dans le repère \( (O ; \vec{i} ; \vec{j}) \).
\end{enumerate}

\exo

\textbf{Partie A :} Soit la fonction \( f \) définie sur \( \mathbb{R} \) par :
\[
f(x) = \frac{e^x}{e^x + 1} - \ln(e^x + 1).
\]

On note \( (C) \) la représentation graphique de la fonction \( f \) dans un plan muni d’un repère orthonormal \( (O ; \vec{i} ; \vec{j}) \) \textbf{(unité 2 cm)}.

\begin{enumerate}
    \item Étudier les variations de \( f \).
    \item Montrer que \( \lim\limits_{x \to +\infty} \left[ f(x) - 1 + x \right] = 0 \).\\
    Que peut-on en déduire pour \( (C_f) \) ?
    \item Construire \( (C_f) \).
    \item Montrer que \( f \) réalise une bijection de \( ]-\infty ; +\infty[ \) sur \( ]-\infty ; 0[ \).
\end{enumerate}
\textbf{Partie B :} Soit \( g \) la fonction définie par :
\[
f(x) = e^x \ln(1 + e^x)
\]
On note \( (C_g) \) sa courbe représentative.

\begin{enumerate}
    \item Montrer que \( g \) est dérivable sur \( \mathbb{R} \).

    \item Montrer que, pour tout réel \( x \),\\
    \( f'(x) = e^{-x} f(x) \).

    \item Montrer que \( \lim\limits_{x \to +\infty} g(x) = 0 \) et \( \lim\limits_{x \to -\infty} g(x) = 1 \).

    \item En déduire la nature des branches infinies.

    \item Dresser le tableau de variations de \( g \).

    \item Construire \( (C_g) \) dans le repère précédent.
    
    \setcounter{enumi}{6}
    \item 
    \begin{enumerate}
        \item Montrer que \( \dfrac{1}{e^x + 1} = \dfrac{e^{-x}}{e^{-x} + 1} \).
        \item À tout réel \( \beta \), on associe le réel
        \[
        I(\beta) = \int_0^{\beta} g(x)\,dx.
        \]
        Justifier l’existence de \( I(\beta) \).
        \item Calculer \( I(\beta) \) à l’aide d’une intégration par parties.
        \item Calculer \( \lim\limits_{\beta \to +\infty} I(\beta) \).
    \end{enumerate}

\end{enumerate}
\textbf{Partie C :} On considère l’équation différentielle :
\[
(E) : y' + y = \dfrac{e^{-x}}{e^{-x} + 1}.
\]

\begin{enumerate}
    \item Vérifier que la fonction \( g \) étudiée dans la \textbf{partie B} est solution de \( (E) \).

    \item Montrer qu’une fonction \( \phi \) est solution de \( (E) \) si et seulement si \( \phi - g \) est solution de l’équation différentielle \( (E_0) : y' + y = 0 \).

    \item Résoudre \( (E_0) \) et en déduire les solutions de \( (E) \).

    \item Déterminer la solution de \( (E) \) qui s’annule en \( \ln 2 \).
\end{enumerate}

\exo
\textbf{vide}
\exo

\textbf{Partie A :} Soit \( h \) la fonction définie sur \( \mathbb{R} \) par :
\[
h(x) = 1 + (1 - x)e^{2 - x}
\]

\begin{enumerate}
    \item Étudier les variations de \( h \) \textbf{(on ne demande pas de calculer les limites aux bornes de \( D_h \))}.
    \item En déduire le signe de \( h(x) \) sur \( \mathbb{R} \).
\end{enumerate}

\textbf{Partie B :} Soit \( f \) la fonction définie sur \( \mathbb{R} \) par :
\[
f(x) = x(1 + e^{2 - x})
\]
et \( (C_f) \) sa courbe représentative dans un repère orthonormé \( (O ; \vec{i} ; \vec{j}) \) \textbf{(unité 2 cm)}.

\begin{enumerate}
    \item 
    \begin{enumerate}
        \item Étudier les limites de \( f \) en \( +\infty \) et en \( -\infty \).
        \item Préciser la nature de la branche infinie en \( -\infty \).
                \item Calculer \( \lim\limits_{x \to +\infty} \left[ f(x) - x \right] \), puis interpréter le résultat obtenu.

        \item Préciser la position de \( (C_f) \) par rapport à la droite \( (\Delta) : y = x \).
    \end{enumerate}
    \item 
    \begin{enumerate} 
            \item Dresser le tableau de variations de \( f \).

            \item Montrer que \( f \) admet une bijection réciproque notée \( f^{-1} \), définie sur \( \mathbb{R} \).

            \item \( f^{-1} \) est-elle dérivable en 4 ?

            \item Étudier la position de \( (C_f) \) par rapport à sa tangente au point d’abscisse 2.

            \item Construire \( (C_f) \) \textbf{(on tracera la tangente au point d’abscisse 2)}.

            \item Construire \( (C_{f^{-1}}) \), la courbe de \( f^{-1} \) dans le repère précédent.
    \end{enumerate}
\end{enumerate}
\textbf{Partie C :} Soit \( \beta \) un réel strictement positif sur \( \mathbb{R} \). La région du plan est délimitée par les droites d’équations \( x = 0 \) et \( x = \beta \), et les courbes d’équations respectives \( y = f(x) \) et \( y = x \). Soit \( A(\beta) \) l’aire de cette région \( R_\beta \), en \( \text{cm}^2 \).

\begin{enumerate}
    \item Calculer \( A(\beta) \) en fonction de \( \beta \).
    \item Déterminer \( \alpha = \lim\limits_{\beta \to +\infty} A(\beta) \).\\
    Interpréter graphiquement le résultat obtenu.
\end{enumerate}

\exo
\textbf{vide}
\exo

\textbf{Partie A :} Soit \( g \) la fonction définie sur \( ]0 ; +\infty[ \) par :
\[
g(x) = 1 + x + \ln x.
\]

\begin{enumerate}
    \item Dresser le tableau de variations de \( g \).
    \item Montrer qu’il existe un unique \( \alpha \) solution de l’équation \( g(x) = 0 \).\\
    Vérifier que \( 0{,}2 < \alpha < 0{,}3 \).
    \item En déduire le signe de \( g \) sur \( ]0 ; +\infty[ \).
    \item Établir la relation : \( \ln \alpha = -1 - \alpha \).
\end{enumerate}

\textbf{Partie B :} On considère la fonction \( f \) définie par :
\[
f(x) = 
\begin{cases}
\dfrac{x \ln x}{1 + x} & \text{si } x > 0 \\
0 & \text{si } x = 0
\end{cases}
\]
et \( (C_f) \) sa courbe représentative dans un repère orthonormé \( (O ; \vec{i} ; \vec{j}) \) \textbf{(unité 5 cm)}.

\begin{enumerate}
    \item Montrer que la fonction \( f \) est continue en 0 puis sur \( [0 ; +\infty[ \).
    \item Étudier la dérivabilité de \( f \) en 0.\\
    Interpréter graphiquement ce résultat.
    \item Déterminer la limite de \( f \) en \( +\infty \).
    \item Montrer que \( \forall x \in ]0 ; +\infty[ \),
    \[
    f'(x) = \dfrac{g(x)}{(1 + x)^2}.
    \]
    En déduire le signe de \( f'(x) \) sur \( ]0 ; +\infty[ \).
    \item Montrer que \( f(\alpha) = -\alpha \).
    \item Dresser le tableau de variations de \( f \).
    \item Construire \( (C_f) \). \textbf{(on prendra \( \alpha = 0{,}3 \)).}
\end{enumerate}
\textbf{Partie C :} Soit \( h \) la restriction de \( f \) à l’intervalle \( I = [1 ; +\infty[ \).

\begin{enumerate}
    \item Montrer que \( h \) réalise une bijection de \( I \) vers un intervalle \( J \) à préciser.
    \item Soit \( h^{-1} \) la bijection réciproque de \( h \).\\
    Étudier la dérivabilité de \( h^{-1} \) sur \( J \).
    \item Calculer \( h(2) \) et \( \left(h^{-1}\right)'\left( \dfrac{2 \ln 2}{3} \right) \).
    \item Construire \( (C_{h^{-1}}) \), la courbe de \( h^{-1} \) dans le repère \( (O ; \vec{i} ; \vec{j}) \).
\end{enumerate}


\textbf{Partie D :}

\begin{enumerate}
    \item À l’aide d’une intégration par parties, calculer l’intégrale 
    \[
    I = \int_1^e x \ln x\,dx.
    \]

    \item Montrer que pour tout \( x \in [1 ; e] \),
    \[
    \dfrac{x \ln x}{e + 1} \leq f(x) \leq \dfrac{x \ln x}{2}.
    \]

    \item En déduire que :
    \[
    \dfrac{e^2 + 1}{4(e + 1)} \leq \int_1^e f(x)\,dx \leq \dfrac{e^2 + 1}{8}.
    \]
\end{enumerate}
\exo
\textbf{vide}
\exo

\textbf{Partie A :} Soit \( f \) la fonction numérique définie par :
\[
f(x) = 
\begin{cases}
x + 2 + \ln\left| \dfrac{x - 1}{x + 1} \right| & \text{si } x < 0 \\
(2 + x)e^{-x} & \text{si } x \geq 0
\end{cases}
\]
et \( (C_f) \) sa courbe représentative dans un repère orthonormé \( (O ; \vec{i} ; \vec{j}) \) \textbf{(unité 1 cm)}.

\begin{enumerate}
    \item Montrer que \( f \) est définie sur \( \mathbb{R} \setminus \{-1\} \).
    
    \item
    \begin{enumerate}
        \item Calculer les limites aux bornes du domaine de définition de \( f \).\\
        Préciser les asymptotes parallèles aux axes.

        \item Calculer \( \lim\limits_{x \to +\infty} \left[ f(x) - (x - 2) \right] \).\\
        Interpréter graphiquement le résultat.
    \end{enumerate}
    
    \item
    \begin{enumerate}
        \item Étudier la continuité de \( f \) en 0.

        \item Démontrer que : 
        \[
        \lim\limits_{x \to 0} \dfrac{e^{-x} - 1}{x} = -1
        \quad \text{et} \quad
        \lim\limits_{x \to 0} \dfrac{\ln(1 - x)}{x} = -1.
        \]

        \item En déduire que \( f \) est dérivable à gauche et à droite en 0.\\
        \( f \) est-elle dérivable en 0 ?
    \end{enumerate}

    \item Calculer \( f'(x) \) pour :
    \begin{enumerate}
        \item \( x \in ]0 ; +\infty[ \).
        \item \( x \in ]-\infty ; -1[ \cup ]-1 ; 0[ \).
    \end{enumerate}

    \item Étudier le signe de \( f'(x) \) pour :
    \begin{enumerate}
        \item \( x \in ]0 ; +\infty[ \).
        \item \( x \in ]-\infty ; -1[ \cup ]-1 ; 0[ \).
    \end{enumerate}

    \item Dresser le tableau de variations de \( f \).
    \item Montrer que l’équation \( f(x) = 0 \) admet une unique solution \( \alpha \) appartenant à \( ]-3 ; -2[ \).
    
    \item Tracer \( (C_f) \). On mettra en évidence l’allure de \( (C_f) \) au point d’abscisse 0 et les droites asymptotes.
\end{enumerate}

\textbf{Partie B :} Soit \( g \) la restriction de \( f \) à l’intervalle \( ]-\infty ; -1[ \).

\begin{enumerate}
    \item Montrer que \( g \) définit une bijection de \( ]-\infty ; -1[ \) sur un intervalle \( J \) à préciser.
    
    \item On note \( g^{-1} \) sa bijection réciproque.
    \begin{enumerate}
        \item Calculer \( g(-2) \). Montrer que \( g^{-1} \) est dérivable en \( \ln 3 \).
        \item Calculer \( (g^{-1})'(\ln 3) \).
        \item Représenter la courbe de \( g^{-1} \) dans le repère précédent.
    \end{enumerate}
\end{enumerate}

\textbf{Partie C :} Soit \( A \) l’aire de la région du plan délimitée par les droites d’équations respectives \( x = -2 \), \( x = -3 \) et \( y = x + 2 \) et la courbe de \( f \).
\exo
\textbf{vide}
\exo

\begin{enumerate}
    \item 
    \begin{enumerate}
        \item Étudier les variations de la fonction \( f \) définie sur \( ]-1 ; +\infty[ \) par : \( f(x) = 2 \ln(x + 1) \).\\
        Tracer sa courbe représentative \( (C) \) et \( (T) \) dans le repère orthonormal \( (O ; \vec{i} ; \vec{j}) \), unité 2 cm.

        \item Démontrer que sur \( [2 ; +\infty[ \), la fonction \( l \), définie par \( l(x) = f(x) - x \), est bijective et que l’équation \( l(x) = 0 \) admet une solution unique \( \lambda \).
    \end{enumerate}

    \item On considère la suite \( (U_n)_{n \in \mathbb{N}} \) définie par :
    \(
    \begin{cases}
    U_0 = 5 \\
    U_{n+1} = 2 \ln(1 + U_n)
    \end{cases}
    \)

    \begin{enumerate}
        \item Sans faire de calcul, représenter les quatre premiers termes de la suite sur le graphique.

        \item Démontrer par récurrence que pour tout \( n \), \( U_n \geq 2 \).

        \item Montrer que pour tout \( x \in [2 ; +\infty[ \), on a :
        \(
        |f'(x)| \leq \dfrac{2}{3}.
        \)

        \item En déduire que pour tout \( n \), on a :
        \(
        |U_{n+1} - \lambda| \leq \dfrac{2}{3} |U_n - \lambda|,
        \)
        \(
        |U_n - \lambda| \leq 2 \left( \dfrac{2}{3} \right)^n,
        \)
        et que \( (U_n) \) converge vers \( \lambda \).

        \item Déterminer le plus petit entier naturel \( p \) tel que :
        \(
        |U_p - \lambda| \leq 10^{-2}.
        \)
        Que représente \( U_p \) pour \( \lambda \) ?
    \end{enumerate}
\end{enumerate}

\exo

\textbf{Partie A :} Soit $g$ la fonction définie sur $]0;+\infty[$ par 
\( g(x)= -\dfrac{2}{1+x^2}+\ln\left(1+\dfrac{1}{x^2}\right). \)

\begin{enumerate}
\item Étudier les variations de $g$ et dresser le tableau de variations de $g$.

\item Montrer que l’équation $g(x)=0$ admet une unique solution $\alpha$ et que :
\( 0,3 \leq \alpha \leq 0,4. \)

\item En déduire le signe de $g$ sur $]0;+\infty[$.
\end{enumerate}

\textbf{Partie B :} Soit la fonction $f$ définie par :

\(
f(x)=
\begin{cases}
x\ln\left(1+\dfrac{1}{x^2}\right) & \text{si } x>0 \\[0.3cm]
-(x+1)e^{-\frac{1}{x}} & \text{si } x<0 \\[0.2cm]
f(0)=0
\end{cases}
\)

On désigne par $(C_f)$ la courbe représentative de $f$ dans un repère orthonormal $(O;\vec{i},\vec{j})$.

\begin{enumerate}
\item Déterminer l’ensemble de définition de $f$ et calculer les limites aux bornes de celui-ci.

\item 
\begin{enumerate}
\item Étudier la continuité puis la dérivabilité de $f$ en $0$ à droite.
\item Interpréter graphiquement les résultats de la question 2.a).
\end{enumerate}

\item Étudier les branches infinies de $(C_f)$.

\item 
\begin{enumerate}
\item Montrer que pour tout $x$ appartenant à $]0;+\infty[$, \( f'(x)=g(x). \)
\item Étudier le sens de variations de $f$ et dresser le tableau de variations de $f$.
\item Montrer que \( f(\alpha)=\dfrac{2\alpha}{1+\alpha^2}. \)
\end{enumerate}

\item Tracer $(C_f)$. (On prendra $\alpha=0,5$, unité 2 cm). On précisera l’équation de la tangente au point d’abscisse $-1$.

\item Soit $I=]-\infty;0[$. Montrer que $f$ est une bijection de $I$ sur $]-\infty;0[$ et préciser l’intervalle $J$.

\item Déterminer l’aire $A(\lambda)$ du domaine compris entre la courbe $(C_f)$, les axes et les droites d’équation $x=1$ et $x=\lambda$ (\( \lambda>0 \)).  
On procédera par intégration par parties.
\end{enumerate}

\exo\

\textbf{Partie A :}

\begin{enumerate}
\item Étudier sur $\mathbb{R}$ le signe de \( 4e^{2x}-5e^x+1 \).

\item Soit $\phi$ la fonction définie par :  
\( \phi(x)=2\sqrt{x}\ln x+2. \)

\begin{enumerate}
\item Déterminer le domaine de définition $D_\phi$ et calculer ses limites aux bornes de $D_\phi$.
\item Étudier ses variations et dresser son tableau de variations.
\item En déduire son signe.
\end{enumerate}
\end{enumerate}

\textbf{Partie B :} Soit la fonction $f$ définie par :
\(
f(x)=
\begin{cases}
x+\dfrac{e^x}{2e^x-1} & \text{si } x\leq 0 \\[0.3cm]
1-x+\sqrt{x}\ln x & \text{si } x>0
\end{cases}
\)

On désigne par $(C_f)$ la courbe représentative de $f$ dans un repère orthonormal $(O;\vec{i},\vec{j})$ (unité 2 cm).

\begin{enumerate}
\item 
\begin{enumerate}
\item Déterminer $D_f$ le domaine de définition de $f$.
\item Calculer les limites aux bornes de $D_f$ et étudier les branches infinies de $(C_f)$.
\item Étudier la position relative de $(C_f)$ par rapport à l’asymptote non parallèle aux axes dans $]-\infty;0[$.
\end{enumerate}

\item 
\begin{enumerate}
\item Étudier la continuité de $f$ en $0$.
\item Étudier la dérivabilité de $f$ en $0$ et interpréter graphiquement les résultats.
\end{enumerate}

\item Déterminer la dérivée de $f$ et dresser le tableau de variations de $f$.

\item Construire dans le repère les asymptotes, la courbe $(C_f)$ et les demi-tangentes.  
On remarquera que \( f(1)=0 \) et \( f'(1)=0 \).

\item Calculer en cm$^2$ l’aire du domaine délimité par $(C_f)$, la droite d’équation $y=x$ et les droites d’équations $x=-\ln 8$ et $x=-\ln 4$.
\end{enumerate}
\exo

Soit $f$ la fonction numérique à variable réelle définie par :
\(
f(x)=
\begin{cases}
x-1+e^{x-1} & \text{si } x\leq 1 \\[0.3cm]
x+\ln\left(\dfrac{1+x}{2x}\right) & \text{si } x>1
\end{cases}
\)

On désigne par $(C_f)$ la courbe représentative de $f$ dans un repère orthonormal $(O;\vec{i},\vec{j})$ unité 1 cm.

\begin{enumerate}
\item Vérifier que $f$ est définie sur $\mathbb{R}$.

\item Étudier la continuité et la dérivabilité de $f$ en $1$ et interpréter graphiquement les résultats.

\item 
\begin{enumerate}
\item Calculer les limites aux bornes de son domaine de définition.
\item Démontrer que la droite $(D_1)$ : \( y=x-1 \) est une asymptote oblique de $(C_f)$ au voisinage de $-\infty$.
\item Démontrer que la droite $(D_2)$ : \( y=x-\ln 2 \) est une asymptote oblique de $(C_f)$ au voisinage de $+\infty$.
\item Étudier les positions relatives de $(D_1)$ et $(D_2)$ respectivement, par rapport à $(C_f)$.
\end{enumerate}

\item Étudier les variations de $f$ puis dresser son tableau de variations.

\item Tracer $(C_f)$ et les droites remarquables.

\item Établir que $f$ réalise une bijection de $\mathbb{R}$ sur un intervalle $J$ à préciser.
\item Tracer $(C_{f^{-1}})$ la courbe de $f^{-1}$, réciproque de $f$, dans le repère précédent.
\begin{enumerate}
\item Calculer en cm$^2$ l’aire du domaine délimité par $(C_f)$, la première bissectrice et les droites d’équations respectives \( x=1 \) et \( x=3 \).
\item En déduire l’aire délimitée par $(C_f)$, $(C_{f^{-1}})$ et les droites d’équations respectives \( x=1 \) et \( x=3 \).  
Justifier votre réponse.
\end{enumerate}
\end{enumerate}
\exo

\textbf{I.} Soit la fonction $f$ définie sur $\mathbb{R}$ par

\( f(x)=\dfrac{3(x-1)^3}{3x^2+1}. \)

\begin{enumerate}
\item Calculer les limites de $f$ aux bornes de son ensemble de définition.

\item Déterminer la dérivée de $f$, étudier son signe et dresser le tableau de variations de $f$.

\item Montrer que l’équation \( f(x)=1 \) admet une solution et une seule $\alpha \in \mathbb{R}$.  
En déduire que \( 3<\alpha<4 \).
\end{enumerate}

\textbf{II.} Soit la fonction $g$ définie par

\( g(x)=\dfrac{3(\ln|x|-1)^3}{3\ln^2|x|+1}. \)
\begin{enumerate}
\item 
\begin{enumerate}
\item Montrer que $g$ est définie sur $\mathbb{R}^*$.
\item Démontrer que $g$ est la composée de la fonction $f$ et d’une fonction $h$ à préciser.
\item Étudier la parité de $g$.
\item On note $D_E=]0;+\infty[$.  
Soit $k$ la restriction de $g$ à $D_E$.  
Calculer les limites de $k$ aux bornes de $D_E$.  
Étudier les branches infinies.
\end{enumerate}

\item 
\begin{enumerate}
\item En utilisant les questions I) et II.1.b), calculer \( k'(x) \) et étudier les variations de $k$ sur $D_E$.  
Dresser le tableau de variations de $k$ sur $D_E$.
\item Déterminer le point d’intersection de la courbe de $k$ avec l’axe des abscisses et préciser le signe de $k$.
\end{enumerate}

\item 
\begin{enumerate}
\item Montrer que $k$ réalise une bijection de $]0;+\infty[$ sur un intervalle $J$ à préciser.
\item Construire les courbes $(C_k)$ et $(C_{k^{-1}})$.  
$(C_{k^{-1}})$ est la courbe représentative de la bijection réciproque $k^{-1}$ de $k$ dans un repère orthonormé ; unité 1 cm.  
Tracer la courbe de $g$ dans le repère précédent.
\end{enumerate}
\end{enumerate}
\exo

\textbf{Partie A :} Soit $f$ la fonction définie par

\(
f(x)=
\begin{cases}
\left( \dfrac{x^2+x+1}{x}\right) e^{-\dfrac{1}{x}} & \text{si } x>0 \\[0.3cm]
\dfrac{x}{x+1}+\ln(x+1) & \text{si } -1<x\leq 0
\end{cases}
\)

et $(C_f)$ sa courbe représentative dans un repère orthonormé $(O;\vec{i},\vec{j})$ (unité 2 cm).

\begin{enumerate}
\item Montrer que l’ensemble de définition de $f$ est :  
\( D_f=]-1;+\infty[. \)

\item 
\begin{enumerate}
\item Montrer les égalités suivantes :  
\( \lim_{x\to -1^+} f(x)=-\infty \) et \( \lim_{x\to +\infty} f(x)=1. \)
\item En déduire les asymptotes de $(C_f)$.
\end{enumerate}

\item 
\begin{enumerate}
\item Montrer que \( \lim_{x\to 0^+} \dfrac{1}{x}e^{-\frac{1}{x}}=0. \)
\item Étudier la continuité de $f$ en $0$.
\item En posant \( h=\dfrac{1}{x} \), étudier  
\( \lim_{x\to 0^+} \dfrac{f(x)-f(0)}{x}. \)
\item Étudier \( \lim_{x\to 0^-} \dfrac{f(x)-f(0)}{x}. \)
\item $f$ est-elle dérivable en $0$ ? Interpréter graphiquement les résultats.
\end{enumerate}

\item 
\begin{enumerate}
\item Montrer que :  

\( \forall x\in ]0;+\infty[,  f'(x)=\left(\dfrac{1-x}{x^4}\right)e^{-\frac{1}{x}}. \)
\item Montrer que

\( \forall x\in ]-1;0[, \quad f'(x)=\dfrac{x+2}{(x+1)^2}. \)
\end{enumerate}

\item Dresser le tableau de variations de $f$.

\item Construire $(C_f)$ dans le repère $(O;\vec{i},\vec{j})$.
\end{enumerate}

\textbf{Partie B :} Soit $\alpha$ un réel tel que : \( -1<\alpha<0. \)

\begin{enumerate}
\item Montrer que \( \forall x\in ]-1;0[, \quad \dfrac{x}{1+x}=1-\dfrac{1}{x+1}. \)

\item En utilisant une intégration par parties, déterminer que : 
 
\( \int_{\alpha}^{0} \ln(1+x)\,dx=-\alpha\ln(1+\alpha)-\ln(1+\alpha). \)

\item En déduire l’aire $A(\alpha)$ du domaine plan délimité par la courbe $(C_f)$, l’axe des abscisses et les droites d’équations \( x=\alpha \) et \( x=0. \)

\item Calculer \( \lim_{\alpha\to -1^+} A(\alpha). \)  
Interpréter géométriquement le résultat.
\end{enumerate}
\exo

\textbf{Partie A :} On considère la fonction $g$ définie par  
\( g(x)=1-x^2-\ln x. \)

\begin{enumerate}
\item Calculer \( g(1). \)

\item Étudier les variations de $g$.

\item Déduire de cette étude le signe de $g(x)$ pour tout $x$ de l’ensemble de définition $D_g$ de la fonction $g$.
\end{enumerate}

\textbf{Partie B :} On considère la fonction $f$ définie par  
\( f(x)=2-x+\dfrac{\ln x}{x}. \)

\begin{enumerate}
\item Déterminer l’ensemble de définition $D_f$ de $f$.

\item Étudier les variations de $f$ et dresser son tableau de variations.

\item 
\begin{enumerate}
\item Montrer que la courbe représentative $(C_f)$ de $f$ admet comme asymptote la droite $(\Delta)$ d’équation : \( y=2-x. \)
\item Étudier la position relative de $(C_f)$ par rapport à $(\Delta)$.
\end{enumerate}

\item Déterminer les coordonnées du point $A$ de $(C_f)$ où la tangente est parallèle à $(\Delta)$.

\item Représenter dans un repère orthonormé $(O;\vec{i},\vec{j})$ (unité graphique 2 cm) la courbe $(C_f)$ et les droites asymptotes.
\end{enumerate}
\exo

\textbf{Partie A :} Soit $g$ la fonction définie sur $[0;+\infty[$ par  
\( g(x)=\dfrac{2x^2}{x^2+1}-\ln(x^2+1). \)

\begin{enumerate}
\item Déterminer \( \lim_{x\to +\infty} g(x). \)

\item Calculer la dérivée de $g$ et donner le tableau de variations de $g$.

\item Montrer que sur $[0;+\infty[$ l’équation \( g(x)=0 \) admet une unique solution $\alpha$ et que  
\( 1,9<\alpha<2. \)

\item Préciser le signe de $g$ sur $[0;+\infty[$.
\end{enumerate}

\textbf{Partie B :} Soit $f$ la fonction définie sur $\mathbb{R}$ par :
\(
f(x)=
\begin{cases}
xe^{\frac{1}{x}} & \text{si } x<0 \\[0.3cm]
\dfrac{\ln(1+x^2)}{x} & \text{si } x>0 \\[0.2cm]
f(0)=0
\end{cases}
\)

et $(C_f)$ sa courbe représentative dans un repère orthonormé $(O;\vec{i},\vec{j})$. Unité 2 cm.

\begin{enumerate}
\item Étudier la continuité et la dérivabilité de $f$ en $0$. Interpréter graphiquement les résultats.

\item Montrer que la droite d’équation \( y=x+1 \) est asymptote à $(C_f)$ en $-1$.

\item Calculer la limite de $f(x)$ quand $x$ tend vers $+\infty$. Interpréter graphiquement le résultat.

\item Calculer \( f'(x) \) dans chacun des intervalles où $f$ est dérivable et donner une relation liant \( f'(x) \) et $g(x)$ pour $x>0$.

\item Établir que \( f(\alpha)=\dfrac{2\alpha}{\alpha^2+1} \) et en déduire un encadrement de $f(\alpha)$.
\item Donner le tableau de variations de $f$ et tracer la courbe $(C_f)$.  
(On prendra \( \alpha=1,95 \) et \( f(\alpha)=0,85 \).)
\end{enumerate}

\textbf{Partie C :} Soit $h$ la restriction de $f$ à l’intervalle $]-\infty;0[$.

\begin{enumerate}
\item Montrer que $h$ est une bijection de $]-\infty;0[$ vers un intervalle $J$ à préciser.

\item Construire $(C_{h^{-1}})$, la courbe de $h^{-1}$ dans le repère $(O;\vec{i},\vec{j})$.
\end{enumerate}
\exo

Les résultats de la partie A seront utilisés dans la partie B.

\textbf{Partie A :}

\begin{enumerate}
\item Montrer que \( \lim_{x\to 0} \dfrac{e^x-x-1}{x}=0. \)

\item Soit $k$ la fonction définie sur $]0;+\infty[$ par  

\( k(x)=x(1-\ln x). \)

\begin{enumerate}
\item $k$ est-elle continue sur $]0;+\infty[$ ? Justifier la réponse.
\item Soit $K$ la fonction définie sur $]0;+\infty[$ par :  
\( K(x)=\dfrac{3}{4}x^2-\dfrac{1}{2}x^2\ln x. \)  
Vérifier que $K$ est une primitive de $k$ dans $]0;+\infty[$.
\end{enumerate}
\end{enumerate}

\textbf{Partie B :} Soit $f$ la fonction définie sur $\mathbb{R}$ par
\(
f(x)=
\begin{cases}
e^x-x-1 & \text{si } x\leq 0 \\[0.3cm]
x\ln x & \text{si } x>0
\end{cases}
\)

Le plan est rapporté à un repère orthonormé $(O;\vec{i},\vec{j})$ (unité 2 cm).

\begin{enumerate}
\item Déterminer $D_f$ puis calculer les limites aux bornes de $D_f$.

\item 
\begin{enumerate}
\item Étudier la continuité de $f$ en $0$.
\item Étudier la dérivabilité de $f$ en $0$. Interpréter géométriquement les résultats.
\end{enumerate}

\item Donner les domaines de continuité et de dérivabilité de $f$.

\item Calculer la dérivée de $f$ sur son domaine d’existence et étudier son signe.

\item Dresser le tableau de variations de $f$.

\item Montrer que la droite d’équation $(\Delta)$ : \( y=-x-1 \) est une asymptote à la courbe $(C_f)$ de $f$ quand $x$ tend vers $-\infty$.

\item Représenter graphiquement la courbe $(C_f)$ dans le repère $(O;\vec{i},\vec{j})$.

\item Soit $h$ la restriction de $f$ à l’intervalle  

\( I=\left[\dfrac{1}{e};+\infty\right[. \)

\begin{enumerate}
\item Montrer que $h$ est une bijection de $I$ vers un intervalle $J$ à préciser.
\item Construire $(C_{h^{-1}})$, la courbe de $h^{-1}$ dans le repère $(O;\vec{i},\vec{j})$.
\end{enumerate}
\item Soit $A_1$ l’aire du domaine du plan délimité par  
\( x=\dfrac{1}{e}, \; x=e \), la courbe $(C_f)$ et la droite $(D)$ d’équation \( y=x \).

\begin{enumerate}
\item Calculer $A_1$.
\item En déduire l’aire du domaine du plan délimité par les droites d’équations respectives  
\( x=-\dfrac{1}{e}, \; x=\dfrac{1}{e} \), la droite $(D)$ et la courbe $(C_{h^{-1}})$.
\end{enumerate}
\end{enumerate}
\exo

\textbf{Partie A :} Soit \( g(x)=(2x-1)e^x+1. \)

\begin{enumerate}
\item Étudier les variations de $g$.

\item En déduire le signe de $g(x)$ pour tout \( x\in ]0;+\infty[. \)
\end{enumerate}

\textbf{Partie B :} Soit $f$ la fonction définie par

\(
f(x)=
\begin{cases}
\dfrac{x}{\ln(-x)} & \text{si } x<0 \\[0.3cm]
\dfrac{e^x-1}{\sqrt{x}} & \text{si } x>0 \\[0.2cm]
f(0)=0
\end{cases}
\)

et $(C_f)$ sa courbe représentative dans un repère orthonormé $(O;\vec{i},\vec{j})$. Unité 2 cm.

\begin{enumerate}
\item Vérifier que $f$ existe pour tout réel $x\neq -1$.

\item Étudier la continuité et la dérivabilité de $f$ en $0$. Interpréter graphiquement les résultats.

\item 
\begin{enumerate}
\item Trouver les limites de $f$ aux bornes de $D_f$.
\item Étudier les branches infinies de $(C_f)$.
\end{enumerate}

\item Montrer que pour tout $x>0$,  
\( f'(x)=\dfrac{g(x)}{2x\sqrt{x}}. \)

\item Calculer $f'(x)$ pour $x<0$ et $x\neq -1$ et dresser le tableau de variations de $f$.

\item Tracer la courbe $(C_f)$.
\end{enumerate}

\textbf{Partie C :} Soit $h$ la restriction de $f$ à l’intervalle $[0;+\infty[$.

\begin{enumerate}
\item Montrer que $h$ est une bijection de $[0;+\infty[$ vers un intervalle $J$ à déterminer.

\item Construire $(C_{h^{-1}})$, la courbe $h^{-1}$ dans le repère $(O;\vec{i},\vec{j})$.
\end{enumerate}
\exo

\textbf{Partie A :} Soit $g$ la fonction définie dans $]0;+\infty[$ par  
\( g(x)=\dfrac{x}{x-1}-\ln|x-1|. \)

\begin{enumerate}
\item 
\begin{enumerate}
\item Déterminer l’ensemble de définition $D_g$ de $g$.
\item Calculer les limites aux bornes de $D_g$. (Pour la limite en $1$, on pourra poser $h=x-1$.)
\end{enumerate}

\item Déterminer $g'$ la dérivée de $g$ et dresser le tableau de variations de $g$.

\item Montrer que l’équation \( g(x)=0 \) admet une unique solution $\alpha$ telle que \( 4<\alpha<5 \).

\item Déduire de l’étude précédente le signe de $g$ sur $D_g$.
\end{enumerate}

\textbf{Partie B :} Soit $f$ la fonction définie par

\(
f(x)=
\begin{cases}
\dfrac{\ln|x-1|}{x} & \text{si } x>0 \\[0.3cm]
\dfrac{6e^x}{e^{2x}+3e^x+2} & \text{si } x\leq 0
\end{cases}
\)

et $(C_f)$ sa courbe représentative dans un repère orthonormé $(O;\vec{i},\vec{j})$ (unité 2 cm).

\begin{enumerate}
\item 
\begin{enumerate}
\item Vérifier que $f$ est bien définie sur $\mathbb{R}\setminus\{1\}$ et calculer les limites aux bornes de son ensemble de définition.
\item Préciser les droites asymptotes à $(C_f)$.
\end{enumerate}

\item 
\begin{enumerate}
\item Étudier la continuité de $f$ en $0$.
\item On admet que  
\( \lim_{x\to 0^+}\dfrac{\ln(1-x)+x}{x^2}=\dfrac{1}{2}. \)  

Montrer que  
\( \lim_{x\to 0^-}\dfrac{f(x)-f(0)}{x}=-\dfrac{1}{6}. \)  

Donner une interprétation graphique de ces résultats.
\end{enumerate}

\item 
\begin{enumerate}
\item Montrer que \( f(\alpha)=\dfrac{1}{\alpha+1}. \)
\item Calculer $f'(x)$ sur les intervalles où $f$ est dérivable puis dresser le tableau de variations de $f$.
\end{enumerate}

\item Construire $(C_f)$. On pourra prendre $\alpha=4,5$.  
On placera les points d’abscisses $-1; 0; 2$ et $5$.

\item 
\begin{enumerate}
\item Déterminer les réels $a$ et $b$ tels que pour tout $x\in \mathbb{R}\setminus\{-2;-1\}$, on ait :  

\( -\dfrac{6x}{x^2+3x+2}=\dfrac{a}{x+2}+\dfrac{b}{x+1}. \)

\item En déduire que :  

\( -\dfrac{e^x}{e^{2x}+3e^x+2}=\dfrac{-12e^{-x}}{1+2e^{-x}}+\dfrac{6e^{-x}}{e^x+1}. \)

\item Calculer l’aire du domaine plan limité par $(C_f)$, l’axe des abscisses et les droites d’équations respectives $x=-\ln 2$ et $x=0$.
\end{enumerate}
\end{enumerate}
\exo
Vide
\exo

\textbf{Partie A :}

\begin{enumerate}
\item En utilisant une intégration par parties, calculer pour tout réel $\alpha$,  
\( I(\alpha)=\int_{0}^{\alpha} e^t(t+1)\,dt. \)  
En déduire $I(\alpha)$.

\item Soit $k$ une fonction dérivable sur $\mathbb{R}$.  
Considérons la fonction $h$ telle que  
\( h(x)=k(x)e^{-x}, \quad \forall x\in\mathbb{R}. \)

On se propose de déterminer la fonction $h$ de la façon à ce qu’elle vérifie les conditions suivantes :

\(
\begin{cases}
h'(x)+h(x)=x+2 \\
h(0)=2
\end{cases}
\)

\begin{enumerate}
\item Vérifier que \( k'(x)=(x+2)e^{-x}. \)
\item En déduire $k$ puis $h$.
\end{enumerate}
\end{enumerate}

\textbf{Partie B}

\textbf{I.} Soit $g$ la fonction définie sur $\mathbb{R}$ par 
 
\( g(x)=x+1+e^{-x}. \)

\begin{enumerate}
\item Étudier les variations de $g$.
\item En déduire que $g(x)$ est strictement positif.
\end{enumerate}

\textbf{II.} Soit $f$ la fonction définie par  

\( f(x)=\ln(x+1+e^{-x}) \) sa courbe représentative dans le plan rapporté à un repère orthonormé $(O;\vec{i},\vec{j})$.

\begin{enumerate}
\item Étudier les variations de $f$ puis dresser son tableau de variations.

\item Pour tout $x$ strictement positif, on note $M$, le point de la courbe de la fonction logarithme népérien d’abscisse $x$ et $N$ le point de $(C_f)$ de même abscisse.

\begin{enumerate}
\item Démontrer que  
\( 0<\overline{MN}<\ln\left(\dfrac{x+2}{x}\right). \)

\item Quelle est la limite de $\overline{MN}$ quand $x$ tend vers $+\infty$.
\end{enumerate}
\item 
\begin{enumerate}
\item Démontrer que :

 \( \forall x\in\mathbb{R}, \quad f(x)=-x+\ln(xe^x+e^x+1). \)
\item En déduire que $(C_f)$ admet une asymptote oblique $(\Delta)$ au voisinage de $+\infty$ et déterminer la position relative de $(C_f)$ par rapport à $(\Delta)$ pour $x<-1$.
\end{enumerate}
\item Construire $(C_f)$ et $(\Delta)$ dans le repère.
\end{enumerate}
\exo

\textbf{Partie A : Questions de Cours}

\begin{enumerate}
\item Rappeler les limites usuelles suivantes :  
\( \lim_{x\to 0^+} x\ln x \) , \quad
\( \lim_{x\to -\infty} xe^x. \)

\item Soit $f$ une fonction définie sur un intervalle $I$ non vide de $\mathbb{R}$ vers un intervalle $J$ non vide de $\mathbb{R}$.

\begin{enumerate}
\item Quand est-ce que l’on dit que $g$ réalise une bijection de $I$ vers $J$ ?
\item Quand est-ce que l’on dit que $g^{-1}$, la bijection réciproque de $g$, est dérivable sur $J$ ?
\end{enumerate}
\end{enumerate}

\textbf{Partie B :} Soit $f$ la fonction définie par

\(
f(x)=
\begin{cases}
1-x+x\ln(-x) & \text{si } x<0 \\[0.3cm]
(2x+1)e^{-\frac{1}{2}x} & \text{si } x\geq 0
\end{cases}
\)

On désigne par $(C_f)$ sa courbe représentative dans un repère orthonormé $(O;\vec{i},\vec{j})$ (unité 2 cm).

\begin{enumerate}
\item Déterminer $D_f$ le domaine de définition de la fonction $f$.

\item Calculer les limites aux bornes de $D_f$.

\item Étudier la nature de chaque branche infinie.

\item Étudier la continuité et la dérivabilité de $f$ en $0$. Interpréter graphiquement les résultats.

\item Calculer $f'(x)$ puis dresser le tableau de variations de $f$.

\item Montrer que l’équation \( f(x)=0 \) admet une unique solution $\alpha$ dans l’intervalle $]-\infty;-1]$ et que \( -4<\alpha<-3. \)

\item En déduire le nombre de points d’intersection de $(C_f)$ avec l’axe des abscisses.

\item Construire $(C_f)$ dans le repère $(O;\vec{i},\vec{j})$.
\end{enumerate}

\textbf{Partie C :} Soit $g$ la restriction de $f$ à l’intervalle  
\( I=]-\infty;-1]. \)

\begin{enumerate}
\item Montrer que $g$ admet une bijection réciproque $g^{-1}$ dont on déterminera l’ensemble de définition.

\item Construire dans le repère $(O;\vec{i},\vec{j})$ $(C_{g^{-1}})$ la courbe de $g^{-1}$.
\end{enumerate}

\exo
\textbf{Partie A :} Soit $g$ la fonction définie par 
 
\( g(x)=\dfrac{x}{x+1}-2\ln(x+1). \)

\begin{enumerate}
\item Déterminer $D_g$, le domaine de définition de $g$ puis calculer les limites aux bornes de $D_g$.

\item Calculer $g'(x)$, étudier son signe, puis dresser le tableau de variations de $g$.

\item Calculer $g(0)$. Montrer que l’équation \( g(x)=0 \) admet exactement deux solutions dont l’une que l’on désigne par $\alpha \in ]-0,72;-0,71[$.

\item Déterminer le signe de $g(x)$.
\end{enumerate}

\textbf{Partie B :} Soit $f$ la fonction définie par

\(
f(x)=
\begin{cases}
\dfrac{x^2}{\ln(x+1)} & \text{si } x>-1 \\[0.3cm]
x+1+e^{-x-1} & \text{si } x\leq -1 \\[0.2cm]
f(0)=0
\end{cases}
\)

et $(C_f)$ sa courbe représentative dans un repère orthonormé $(O;\vec{i},\vec{j})$ (unité 2 cm).

\begin{enumerate}
\item 
\begin{enumerate}
\item Montrer que $D_f=\mathbb{R}$ et calculer les limites aux bornes de $D_f$.
\item Étudier la nature des branches infinies.
\end{enumerate}

\item 
\begin{enumerate}
\item Étudier la continuité de $f$ en $-1$ et en $0$.
\item Étudier la dérivabilité de $f$ en $-1$ et en $0$. Interpréter graphiquement les résultats.
\end{enumerate}

\item 
\begin{enumerate}
\item Montrer que pour tout \( x\in ]-1;+\infty[ \) et \( x\neq 0 \),  
\( f'(x)=-\dfrac{xg(x)}{\ln^2(x+1)} \)  
et calculer $f'(x)$ sur \( ]-\infty;-1[ \).
\item Étudier les variations de $f$ et dresser son tableau de variations.
\end{enumerate}

\item Soit $h$ la restriction de $f$ à $[0;+\infty[$.

\begin{enumerate}
\item Montrer que $h$ réalise une bijection de $[0;+\infty[$ sur un intervalle $J$ à préciser.
\item Donner le sens de variations de $h^{-1}$.
\item Construire $(C_f)$ et $(C_{h^{-1}})$.
\end{enumerate}
\end{enumerate}

\textbf{Partie C :} Soit $m$ la fonction définie sur $]0;+\infty[$ par :  
\( m(x)=\dfrac{\ln(x+1)}{x^2}-\dfrac{1}{x(x+1)}. \)

\begin{enumerate}
\item Déterminer les fonctions $u$ et $v$ telles que :  

\( m(x)=u'(x)v(x)+v'(x)u(x). \)

\item En déduire la fonction $H$ définie sur $]0;+\infty[$ telle que \( H'(x)=m(x) \) puis calculer  
\( \displaystyle \int_{1}^{2} \dfrac{1}{f(x)}\,dx. \)
\end{enumerate}

\exo

\textbf{Partie A :}

Soit $g$ la fonction définie par :  
\( g(x)=x^2-1+\ln x. \)

\begin{enumerate}
\item Étudier les variations de $g$.
\item Calculer $g(1)$ puis en déduire le signe de $g(x)$.
\end{enumerate}

\textbf{Partie B :}

\(
f(x)=
\begin{cases}
\dfrac{-x+1}{e^x} & \text{si } x\leq 1 \\[0.3cm]
\dfrac{-x+1}{2}+\dfrac{\ln x}{2x} & \text{si } x>1
\end{cases}
\)

\begin{enumerate}
\item Montrer que $f$ est définie sur $\mathbb{R}$.
\item Étudier la continuité de $f$ en $1$.
\item Étudier la dérivabilité de $f$ en $1$. Interpréter graphiquement le résultat.
\item Déterminer les limites de $f$ aux bornes de son ensemble de définition.
\item Étudier les branches infinies.
\item 
\begin{enumerate}
\item Déterminer $f'$ et étudier son signe.
\item Dresser le tableau de variations de $f$.
\end{enumerate}
\item Représenter $f$ dans un repère orthonormé $(O;\vec{i},\vec{j})$.
\item 
\begin{enumerate}
\item Établir que $f$ admet une bijection réciproque $f^{-1}$.
\item Représenter $f^{-1}$ dans le repère précédent.
\end{enumerate}
\item Soit $h$ la restriction de $f$ à l’intervalle \( \left[0;\dfrac{1}{2}\right] \).  
Montrer que l’équation \( h(x)=x \) admet une solution unique \( \alpha \in \left[0;\dfrac{1}{2}\right] \).
\item Calculer  
\( I=\int_{0}^{\frac{1}{2}} h(x)\,dx \)  
à l’aide d’une intégrale par parties.
\end{enumerate}

\exo

\textbf{Partie A :} Soit $f$ la fonction définie par :

\(
f(x)=
\begin{cases}
2x+2+\ln(2e^{-x}-1) & \text{si } x\leq 0 \\[0.3cm]
\dfrac{1}{x(1-\ln x)} & \text{si } x>0
\end{cases}
\)

On désigne par $(C_f)$ sa courbe représentative dans un repère orthonormé $(O;\vec{i},\vec{j})$ (unité 2 cm).

\begin{enumerate}
\item 
\begin{enumerate}
\item Montrer que $D_f=\mathbb{R}\setminus\{e\}$ puis calculer les limites aux bornes de $D_f$.
\item Montrer que pour tout $x\leq 0$,  

\( f(x)=x+2+\ln(1-e^x). \)
\item Étudier le signe de $x(1-\ln x)$.
\item On admet que  
\( \lim_{x\to -\infty}\dfrac{\ln(2-e^x)}{x}=-1. \)  
Calculer  
\( \lim_{x\to -\infty}\dfrac{f(x)-f(0)}{x} \)  
et interpréter géométriquement le résultat.
\item Montrer que la droite $(\Delta)$ d’équation

\( y=x+2+\ln 2 \) est une asymptote à $(C_f)$ au voisinage de $-\infty$, puis étudier la position relative de $(C_f)$ et de $(\Delta)$.
\end{enumerate}

\item 
\begin{enumerate}
\item Étudier les variations de $f$.
\item Dresser le tableau de variations de $f$.
\item Montrer que l’équation \( f(x)=0 \) admet une unique solution $\alpha$ située dans l’intervalle $]-3;-2[$.  
En déduire une valeur approchée de $\alpha$ à $10^{-1}$ près.
\end{enumerate}

\item Tracer les droites asymptotes à $(C_f)$ puis la courbe $(C_f)$.
\end{enumerate}

\textbf{Partie B :} Soit $g$ la restriction de $f$ à l’intervalle  
\( I=]e;+\infty[. \)

\begin{enumerate}
\item Montrer que $g$ réalise une bijection de $I$ vers un intervalle $J$ à préciser.
\item On note $g^{-1}$ la bijection réciproque de $g$.
\begin{enumerate}
\item Dresser le tableau de variations de $g^{-1}$.
\item Comment obtient-on la courbe $(C_{g^{-1}})$ à partir de la courbe $(C_g)$ ?
\end{enumerate}
\end{enumerate}

\textbf{Partie C :} Soit $F$ la fonction définie par  

\( F(x)=\ln|1-\ln x|. \)

\begin{enumerate}
\item Déterminer l’ensemble de définition de $F$.
\item Déterminer la fonction $F'$ dérivée de $F$.
\item En déduire l’aire du domaine délimité par la courbe $(C_f)$, l’axe des abscisses et les droites d’équations respectives $x=3$ et $x=5$.
\end{enumerate}

\exo

\textbf{Partie A :}

On considère la fonction $g$ définie sur $]0;+\infty[$ par :  
\( g(x)=-\dfrac{1}{x+1}+\ln(x+1)-\ln x. \)

\begin{enumerate}
\item Calculer \( \lim_{x\to +\infty} \ln\left(1+\dfrac{1}{x}\right) \) puis  
\( \lim_{x\to +\infty} g(x). \)

\item Étudier les variations de $g$.

\item En déduire que pour tout $x>0$, $g(x)>0$.
\end{enumerate}

\textbf{Partie B :}

On considère la fonction $f$ définie par :
\(
\begin{cases}
f(0)=0 \\
f(x)=x[\ln(x+1)-\ln x] & \text{si } x>0
\end{cases}
\)

\begin{enumerate}
\item Étudier la continuité de $f$ sur $\mathbb{R}_+$.

\item Étudier la dérivabilité de $f$ en $0$. Interpréter géométriquement le résultat.

\item Calculer $f'(x)$ pour $x>0$. Déterminer le signe de $f'(x)$.

\item Montrer que \( \lim_{x\to +\infty} f(x)=1. \)

\item Dresser le tableau de variations de $f$.

\item Soit $(C_f)$ la courbe représentative de $f$ dans le plan rapporté à un repère orthonormé $(O;\vec{i},\vec{j})$, d’unité graphique 4 cm.

\begin{enumerate}
\item Soit $(T)$ la tangente à $(C_f)$ au point d’abscisse $1$.  
Déterminer l’intersection de l’axe des abscisses avec $(T)$.
\item Construire $(T)$ et $(C_f)$.
\end{enumerate}

\item Soit $\alpha$ un réel strictement positif.

\begin{enumerate}
\item Montrer que  
\( \int_{0}^{\alpha} \dfrac{x}{x+1}\,dx=\alpha-\ln(\alpha+1). \)

\item Calculer l’aire du domaine plan délimité par la courbe $(C_f)$, l’axe des abscisses et les droites d’équations $x=\alpha$ et $x=1$.
\end{enumerate}
\end{enumerate}

\textbf{Partie C :}

On considère les suites $(u_n)$ et $(v_n)$ définies pour $n\in\mathbb{N}$ par  
\( u_n=\left(1+\dfrac{1}{n}\right)^n \) et \( v_n=\ln(u_n). \)

\begin{enumerate}
\item À l’aide de la partie B, déterminer le sens de variations de $(v_n)$ et  
\( \lim_{n\to +\infty} v_n. \)

\item En déduire le sens de variations de $(u_n)$ et  
\( \lim_{n\to +\infty} u_n. \)
\end{enumerate}

\exo

\textbf{Partie A :}

\begin{enumerate}
\item Soit l’équation différentielle $(E)$ :  
\( y''+4y'+4y=0. \)  
Déterminer les solutions $h$ de $(E)$ définies sur $\mathbb{R}$.

\item On considère l’équation différentielle $(F)$ : 
 
\( y''+4y'+4y=-4x. \)

\begin{enumerate}
\item Déterminer les réels $a$ et $b$ tels que la fonction $\varphi : x \mapsto ax+b$ soit solution de $(F)$.
\item Montrer qu’une fonction $f$ est solution de $(F)$ si et seulement si $f-\varphi$ est solution de $(E)$.
\item En déduire toutes les solutions de $(F)$.
\item Donner les solutions $f$ de $(F)$ qui vérifient :  
\( f(0)=2 \) et \( f'(0)=-2. \)
\end{enumerate}
\end{enumerate}

\textbf{Partie B :}

On considère la fonction $f$ définie sur l’intervalle  
\( ]-\infty;-1[\cup]0;+\infty[ \) par :

\(
f(x)=
\begin{cases}
\ln\left(\dfrac{x+1}{x}\right) & \text{si } x<-1 \\[0.4cm]
xe^{-2x}+e^{-2x}-x+1 & \text{si } x\geq 0
\end{cases}
\)

\begin{enumerate}
\item 
\begin{enumerate}
\item Calculer $f'$ et $f''$ de la fonction $f$ sur $[0;+\infty[$.
\item Étudier les variations de $f$ puis dresser le tableau de variation de $f$ sur $[0;+\infty[$.
\item En déduire le signe de $f'$ sur $[0;+\infty[$.
\end{enumerate}

\item Étudier les variations de $f$ sur $]-\infty;-1[$.

\item Dresser le tableau de variations de $f$.

\item Montrer que l’équation \( f(x)=0 \) admet une solution unique $\alpha$ et que :  
\( 1<\alpha\leq 2. \)

\item Montrer que la courbe $(C_f)$ admet une asymptote oblique $(D)$ que l’on déterminera, puis étudier la position de $(D)$ par rapport à la courbe $(C_f)$.

\item Construire les asymptotes, puis la courbe $(C_f)$.
\end{enumerate}

\exo

\textbf{Partie A :} Soit la fonction $f$ définie sur $[0;+\infty[$ par :
\(
\begin{cases}
f(x)=x\ln x & \text{si } x>0 \\
f(0)=0
\end{cases}
\)

\begin{enumerate}
\item Étudier la continuité et la dérivabilité de $f$ en $0$. Interpréter graphiquement le résultat obtenu.
\item Étudier la nature de la branche infinie.
\item Étudier les variations de $f$ puis dresser son tableau de variations.
\item On désigne par $(C_f)$ la courbe représentative de la fonction $f$ dans un repère orthonormé $(O;\vec{i},\vec{j})$ tel que \( \|\vec{i}\|=\|\vec{j}\|=4\,\text{cm} \). Tracer $(C_f)$ dans ce repère.
\end{enumerate}

\textbf{Partie B :}

Soit $g$ la fonction numérique d’une variable réelle définie par :  
\( g(x)=\ln x-e^{1-x}+1. \)

\begin{enumerate}
\item Déterminer $D_g$ puis calculer les limites aux bornes de $D_g$.
\item Étudier les variations de $g$ puis dresser son tableau de variations.
\item Calculer $g(1)$ puis en déduire le tableau de signe de $g$ puis le signe de $g(x)$ sur $]0;+\infty[$.
\end{enumerate}

\textbf{Partie C :}

Soit $h$ la fonction numérique d’une variable réelle définie par :
\(
\begin{cases}
h(x)=x\ln x+e^{1-x} & \text{si } x>0 \\
h(0)=e
\end{cases}
\)

\begin{enumerate}
\item 
\begin{enumerate}
\item Montrer que $h$ est continue à droite en $0$.
\item $h$ est-elle dérivable à droite en $0$ ? Interpréter géométriquement le résultat obtenu.
\end{enumerate}

\item Vérifier que $h'(x)=g(x)$. Déduire de la question 3) de la partie B les variations de $h$, dresser son tableau de variations.

\item $(C_h)$ est la courbe représentative de $h$ dans le repère précédent.
\begin{enumerate}
\item Étudier le signe de $h(x)-f(x)$ pour tout $x>0$.
\item En déduire la position relative de $(C_h)$ par rapport à $(C_f)$.
\item Construire $(C_h)$.
\end{enumerate}
\end{enumerate}

\textbf{Partie D :}

Soit $(u_n)$ la suite définie par :
 
\( \forall n\in\mathbb{N}, \quad u_n=\int_n^{n+1} (h(x)-f(x))\,dx. \)

\begin{enumerate}
\item Interpréter géométriquement $u_1$.
\item Calculer $u_n$ en fonction de $n$ et en déduire que la suite $(u_n)$ est une suite géométrique convergente dont on précisera la raison $q$ et le premier terme $u_0$.
\item Soit $(A_n)$ la suite définie par :  

\( A_n=u_0+u_1+\cdots+u_n. \)

\begin{enumerate}
\item Exprimer $(A_n)$ en fonction de $n$.
\item Déterminer la limite de la suite $(A_n)$ quand $n$ tend vers $+\infty$.
\end{enumerate}
\end{enumerate}

\exo

Soit $f$ la fonction définie pour tout réel $x$ par
  
\( f(x)=-x+2+(2x-4)e^{\frac{x}{2}}. \)

On note $(C_f)$ sa courbe représentative dans un repère orthonormé $(O;\vec{i},\vec{j})$ du plan.  
On choisit 2 cm pour unité graphique.

\textbf{Partie A :}

Soit $g$ la fonction numérique définie pour tout réel $x$ par  
\( g(x)=-1+xe^{\frac{x}{2}}. \)

\begin{enumerate}
\item Calculer les limites de $g$ en $+\infty$ et en $-\infty$.

\item Étudier le sens de variations de $g$ puis dresser le tableau de $g$.

\item Démontrer que l’équation \( g(x)=0 \) admet une solution unique $\alpha$ et une seule puis prouver que  
\( 0,70\leq \alpha \leq 0,71. \)

\item En déduire le signe de $g(x)$.
\end{enumerate}

\textbf{Partie B :}

\begin{enumerate}
\item 
\begin{enumerate}
\item Exprimer $f'(x)$ en fonction de $g(x)$.
\item En déduire le sens de variations de $f$.
\item Démontrer que \( f(\alpha)=4-\alpha-\dfrac{4}{\alpha} \), où $\alpha$ est le nombre défini en 3. Partie A.
\item Donner un encadrement de $f(\alpha)$ d’amplitude $0,1$.
\item 
\begin{enumerate}
\item Déterminer la limite de $f(x)$ et de \( \dfrac{f(x)}{x} \) quand $x$ tend vers $+\infty$.
\item Déterminer que $(C_f)$ admet au voisinage de $-\infty$ une asymptote $(D)$ dont on donnera une équation.
\end{enumerate}
\item Dresser le tableau de variations de $f$.
\item À l’aide d’une intégration par parties, calculer pour tout réel $x$ l’intégrale 
 
\( I(x)=\int_{0}^{x}(2t-4)e^{\frac{t}{2}}dt. \)
\item Soit $\lambda$ un réel négatif. Calculer en cm$^2$ l’aire $A$ du domaine constitué des points de coordonnées $(x;y)$ satisfaisant à  
\( \lambda\leq x\leq 0 \) et \( f(x)\leq y\leq 2-x. \)  
Interpréter graphiquement la limite de l’aire $A$ quand $\lambda$ tend vers $-\infty$.
\end{enumerate}
\end{enumerate}

\exo

\textbf{Partie A :}

Soient $g$ et $h$ les fonctions définies respectivement sur $]-\infty;2[$ et sur $]2;+\infty[$ par :  
\( g(x)=(x-1)e^x-1 \)  
et  
\( h(x)=2\ln(x-1)-\dfrac{x^2-2x}{x-1}. \)

\begin{enumerate}
\item 
\begin{enumerate}
\item Calculer $g(2)$ et la limite de $g$ en $-\infty$.
\item Calculer $g'(x)$ et dresser le tableau de variations de $g$.
\item Montrer que l’équation \( g(x)=0 \) admet une unique solution $\alpha$ et que  
\( 1,2<\alpha<1,3. \)
\item Déterminer le signe de $g(x)$ sur $]-\infty;2]$.
\end{enumerate}

\item 
\begin{enumerate}
\item Calculer $h'(x)$ et dresser le tableau de variations de $h$.
\item En déduire le signe de $h(x)$.
\end{enumerate}
\end{enumerate}

\textbf{Partie B :}

On considère la fonction $f$ définie par :

\(
f(x)=
\begin{cases}
(x-2)e^x-x-4 & \text{si } x\leq 2 \\[0.3cm]
\dfrac{x\ln(x-2)}{x-2} & \text{si } x>2
\end{cases}
\)

On note $(C_f)$ sa courbe représentative dans un repère orthonormé $(O;\vec{i},\vec{j})$ tel que  
\( \|\vec{i}\|=\|\vec{j}\|=2\,\text{cm}. \)

\begin{enumerate}
\item 
\begin{enumerate}
\item Montrer que l’ensemble de définition de $f$ est $\mathbb{R}$.
\item Calculer les limites de $f$ aux bornes de son ensemble de définition.
\end{enumerate}

\item 
\begin{enumerate}
\item Montrer que la droite $(\Delta)$ d’équation \( y=-x+4 \) est une asymptote oblique à $(C_f)$ au voisinage de $-\infty$.
\item Étudier la position relative de $(\Delta)$ par rapport à $(C_f)$.
\end{enumerate}

\item Étudier la continuité de $f$ en $2$.

\item Étudier la dérivabilité de $f$ en $2$ et interpréter géométriquement le résultat obtenu.  

On admet que  
\( \lim_{x\to 2^+} \dfrac{x\ln(x-1)-2x+4}{(x-2)^2}=0. \)

\item 
\begin{enumerate}
\item Calculer $f'(x)$ pour \( x\in]-\infty;2[ \) et pour \( x\in]2;+\infty[ \) puis montrer que pour tout  
\( x\in]2;+\infty[ \),  
\( f'(x)=\dfrac{g(x)}{h(x)} \).
\item En déduire des questions 1.d) et 2.b) de la partie A, dresser le tableau de variations de $f$.
\end{enumerate}

\item Construire la courbe $(C_f)$ dans le repère.
\end{enumerate}

\textbf{Partie C :}

Soit $\lambda>0$ et $D_\lambda$ le domaine délimité par la courbe $(C_f)$, la droite $(\Delta)$ et les droites d’équations respectives  
\( x=2 \) et \( x=\lambda. \)

\begin{enumerate}
\item Déterminer en fonction de $\lambda$ l’aire $A(D_\lambda)$ du domaine $D_\lambda$ en cm$^2$.
\item Calculer \( \lim_{\lambda\to +\infty} A(D_\lambda). \)
\end{enumerate}
\end{multicols}

\end{document}